\documentclass[10pt]{article}

\usepackage[margin=0.68in]{geometry}
\usepackage[T1]{fontenc}
\usepackage{lmodern}
\usepackage{microtype}
\usepackage{amsmath,amssymb,booktabs,graphicx}
\usepackage[dvipsnames]{xcolor}
\usepackage{enumitem}
\usepackage{titlesec}
\usepackage{caption}
\usepackage[colorlinks=true,allcolors=MidnightBlue]{hyperref}

\setlength{\parindent}{0pt}
\setlength{\parskip}{3.5pt}
\setlist[itemize]{leftmargin=1.15em,itemsep=1pt,topsep=2pt}
\titlespacing*{\section}{0pt}{5pt}{2pt}
\titleformat{\section}{\large\bfseries}{}{0pt}{}
\captionsetup{font=small,labelfont=bf,skip=3pt}
\newcommand{\Rcal}{\mathcal R}

\begin{document}

\begin{center}
{\LARGE\bfseries A Fertility Shift, Population Decline, and Housing Prices}\par
\vspace{2pt}
{\large A two-state steady-state thought experiment}\par
\vspace{2pt}
{\small August 7, 2026}
\end{center}

\vspace{-3pt}
\textbf{Question.} Can a persistent fall in desired fertility move a closed economy
from one positive steady state to another with fewer households and cheaper housing?
Yes, provided housing costs feed back into fertility. The permanent object is a
downward shift in the \emph{fertility schedule}. Realized fertility is below
replacement during the transition, but the fall in population lowers housing costs
until fertility returns to replacement at a smaller steady state. If realized
fertility instead remained permanently below replacement, a closed population would
converge to zero rather than to a smaller positive level.

\section*{Environment}

\textbf{Demography.} Let $N_t$ denote adult households and let $\Rcal_t$ be the
net reproduction rate, normalized so that $\Rcal_t=1$ replaces the current adult
population. Housing cost $p_t$ reduces desired fertility, while $\theta_t$ shifts
fertility for nonhousing reasons. Around the initial steady state,
\begin{equation}
 \Rcal_t-1=\theta_t-\eta\,\pi_t,
 \qquad
 \dot n_t=\lambda(\Rcal_t-1),
 \label{eq:demography}
\end{equation}
where $n_t\equiv\log(N_t/N_0)$, $\pi_t\equiv\log(p_t/p_0)$,
$\eta>0$ is the fertility response to housing cost, and $\lambda>0$ maps a
reproduction gap into population growth. Equation \eqref{eq:demography} is the
feedback: cheaper housing raises fertility.

\textbf{Housing demand.} With fixed income per household and a constant housing
expenditure share, aggregate expenditure on housing services is proportional to the
number of households:
\begin{equation}
 p_tH_t=A N_t,
 \qquad
 \pi_t=n_t-h_t,
 \label{eq:housing_demand}
\end{equation}
Here $A$ is fixed housing expenditure per household and
$h_t\equiv\log(H_t/H_0)$ is the usable housing stock. Holding the stock fixed, a
smaller population lowers the housing-services price one for one in logs.

\textbf{Housing stock.} Structures depreciate at rate $\delta$, while gross
construction is $X_t=\delta H_0(p_t/p_0)^\xi$ with finite price elasticity
$\xi\geq0$. Log-linearizing $\dot H_t=X_t-\delta H_t$ gives the stock law
\begin{equation}
 \dot h_t=\delta\big[\xi\pi_t-h_t\big],
 \qquad \delta>0.
 \label{eq:housing_stock}
\end{equation}
Thus the stock gradually approaches the long-run supply schedule
$h=\xi\pi$. This is the smallest way to let excess housing disappear through
depreciation rather than instantaneously.

\section*{Steady states and transition}

Initially $\theta_0=0$ and $(n_0,h_0,\pi_0)=(0,0,0)$. At date zero, a permanent
shock sets $\theta_1<0$. The unchanged initial housing cost makes
$\Rcal_0-1=\theta_1<0$, so population begins to fall. Substituting
\eqref{eq:housing_demand} into \eqref{eq:demography}--\eqref{eq:housing_stock}
gives the exact linear system
\begin{equation}
 \begin{pmatrix}\dot n_t\\ \dot h_t\end{pmatrix}
 =
 \underbrace{\begin{pmatrix}
 -a & a\\ \delta\xi & -\delta(1+\xi)
 \end{pmatrix}}_{\mathsf A}
 \begin{pmatrix}n_t\\h_t\end{pmatrix}
 +\begin{pmatrix}\lambda\theta_1\\0\end{pmatrix},
 \qquad a\equiv\lambda\eta>0.
 \label{eq:system}
\end{equation}
The new steady state is pinned down by replacement fertility and housing-stock
balance:
\begin{equation}
 \boxed{
 \pi_1^*=\frac{\theta_1}{\eta}<0,
 \qquad
 h_1^*=\xi\frac{\theta_1}{\eta}\leq0,
 \qquad
 n_1^*=(1+\xi)\frac{\theta_1}{\eta}<0.}
 \label{eq:ss}
\end{equation}
The population decline reduces housing demand; lower prices reduce construction and
allow the stock to depreciate; and the lower price raises fertility back to
$\Rcal^*=1$. A more elastic long-run housing supply makes the stock contract more,
so a larger population decline is required to deliver the price decline that restores
replacement fertility.

The two adjustment roots are
\begin{equation}
 \mu_{\pm}=-\frac{a+\delta(1+\xi)\mp
 \sqrt{[a+\delta(1+\xi)]^2-4a\delta}}{2}<0.
 \label{eq:roots}
\end{equation}
Hence the new steady state is stable. The whole transition is analytical:
$z_t=z_1^*+\exp(\mathsf A t)(z_0-z_1^*)$, where
$z_t=(n_t,h_t)^\top$.

\newpage
\section*{Illustration}

The figure normalizes initial population, housing, and price to one. The permanent
fertility-schedule shift requires a 10 percent housing-price decline to restore
replacement fertility, $\theta_1/\eta=\log(0.90)$. Setting $\xi=1$,
$a=0.02$, and $\delta=0.025$ gives a slow adjustment root of $-0.0081$, or a
half-life of about 86 years. These values illustrate the mechanism; they are not a
calibration or a forecast.

\begin{figure}[h!]
\centering
\begin{minipage}[t]{0.485\textwidth}
\centering
\includegraphics[width=\linewidth]{../output/model/fertility_population_housing_transition_note/population_transition.pdf}
\end{minipage}\hfill
\begin{minipage}[t]{0.485\textwidth}
\centering
\includegraphics[width=\linewidth]{../output/model/fertility_population_housing_transition_note/housing_price_transition.pdf}
\end{minipage}
\caption{After the schedule shifts down, realized fertility is below replacement and
population declines. The housing stock follows slowly, so housing cost falls. The
price decline raises fertility until the net reproduction rate returns to one. In the
illustration, $N_1^*/N_0=0.81$, $H_1^*/H_0=0.90$, and $p_1^*/p_0=0.90$.}
\label{fig:transition}
\end{figure}

\section*{What carries over to the quantitative model}

The toy model isolates three objects already present in the quantitative environment.
First, lower housing costs raise fertility through housing consumption, tenure, and
down-payment margins. Second, lower fertility reduces future adult household entry,
although the full age structure inserts a long cohort delay that is suppressed here.
Third, the level housing-supply curve makes a smaller population cheaper to house. In
a stationary state the asset price satisfies $P^{H,*}=p^*/u$, where $u$ is the
stationary user-cost factor, so the long-run asset-price comparison has the same sign
as \eqref{eq:ss}. A complete asset-price transition would also capitalize future rents
and could jump on impact.

This thought experiment is an \emph{alternative closure}, not a reinterpretation of
the current quota exercise. Under the quota closure, permanently subreplacement
fertility is compatible with a positive population because a fixed outside inflow
$\bar M$ replaces the missing native entrants. A one-state discrete analogue of its
stationary renewal accounting uses the household exit rate $E$, the mature
child-household flow $B_t$, and the retained share $\bar R$:
\begin{equation}
 S_{t+1}=(1-E+\bar R B_t)S_t+\bar M,
 \qquad
 S^*=\frac{\bar M}{E-\bar R B^*}.
 \label{eq:quota}
\end{equation}
With $\bar M=0$ and $0\leq1-E+\bar R B^*<1$, equation \eqref{eq:quota}
instead implies $S_t\rightarrow0$. The closed-economy model above avoids that
outcome because fertility is endogenous and returns to replacement as housing
becomes cheaper.

\section*{Literature map and limits}

The closest Fern\'andez-Villaverde reference is
\href{https://www.nber.org/papers/w29480}{Delventhal, Fern\'andez-Villaverde, and
Guner (2021)}. Their Appendix J analytically characterizes an initial Malthusian
steady state in which fertility equals replacement; fixed land enters production
rather than housing. The paper does not derive a lower-population, lower-house-price
steady state. The likely
``Rognlie'' reference in the handwritten notes is
\href{https://www.brookings.edu/articles/deciphering-the-fall-and-rise-in-the-net-capital-share/}{Rognlie (2015)},
which motivates scarce residential land but has no fertility transition.
\href{https://doi.org/10.1086/427465}{Glaeser and Gyourko (2005)} provide the
durable-housing logic: a negative local-demand shock lowers prices immediately,
while housing and population contract gradually as depreciated units are not rebuilt.
\href{https://www.nber.org/papers/w2794}{Mankiw and Weil (1989)} provide a
partial-equilibrium link from cohort size to housing demand.

The clean sign in Figure \ref{fig:transition} is a mechanism result, not a forecast.
\href{https://doi.org/10.1016/j.econmod.2025.107125}{Lee and Suh (2025)} show in a
full OLG model that aging and a lower equilibrium interest rate can initially raise
house prices even though very low fertility lowers them in the long run. The present
two-state model deliberately holds income, interest rates, taxes, tenure, credit
constraints, and age composition fixed so that the fertility--housing feedback is
visible analytically.

\end{document}
